\documentclass[12pt]{article}
\usepackage[utf8]{inputenc}
\usepackage[T1]{fontenc}
\usepackage[margin=1in]{geometry}
\usepackage{setspace}
\usepackage{xcolor}
\usepackage[colorlinks=true, urlcolor=blue!60!black, linkcolor=blue!60!black]{hyperref}
\setlength{\parindent}{2em}
\setlength{\parskip}{1.25em}
\renewcommand{\baselinestretch}{1.0}
\begin{document}
\begin{center}
{\LARGE \textit{Wolford v. Lopez}, No. 24-1046 (2026)\par}
\vspace{0.8em}
{\large Case Brief\par}
\end{center}
\vspace{1em}

\section{Memory Jogger}
In a 6--3 decision, the Supreme Court struck down Hawaii's post-\textit{Bruen} ``vampire rule'' barring licensed gun owners from carrying handguns onto private property open to the public without the owner's express permission, holding that states cannot flip the common-law implied-license default to hobble public carry (\href{https://www.courtlistener.com/opinion/10880241/wolford-v-lopez/}{CourtListener}; \href{https://www.law.cornell.edu/supremecourt/text/24-1046}{Cornell LII}).

\section{Detailed Case Facts}
Prior to 2022, Hawaii strictly limited public carry of firearms, issuing concealed-carry permits only in ``exceptional case[s]'' where an applicant demonstrated a specific reason to fear injury; between 2000 and 2018, the state reportedly issued only four such permits (\href{https://www.law.cornell.edu/supremecourt/text/24-1046}{Cornell LII}). In \textit{New York State Rifle \& Pistol Ass'n v. Bruen}, 597 U.S. 1 (2022), the U.S. Supreme Court invalidated New York's ``proper cause'' licensing scheme, held that the Second and Fourteenth Amendments protect an individual right to carry handguns outside the home for self-defense, and singled out Hawaii as one of six other jurisdictions with comparably restrictive regimes (\href{https://www.supremecourt.gov/opinions/25pdf/24-1046_nmio.pdf}{Slip Opinion}).

In response to \textit{Bruen}, the Hawaii Legislature enacted Act 52 on June 2, 2023 (2023 Haw. Sess. Laws 113--136), creating a rigorous licensing process under Haw. Rev. Stat. \S{}134-9 requiring background checks, disclosure of medical information, live-fire safety training, and proof of ``essential character or temperament'' (\href{https://www.law.cornell.edu/supremecourt/text/24-1046}{Cornell LII}; \href{https://harvardlawreview.org/print/vol-140/wolford-v-lopez/}{Harvard Law Review}). Even for licensed permit holders, Act 52 imposed two sweeping geographic limits. First, Haw. Rev. Stat. \S{}134-9.1(a) designated numerous ``sensitive places'' where firearms were flatly banned, including government buildings, hospitals, restaurants serving alcohol, stadiums, theaters, libraries, schools, universities, beaches, parks, banks, and public gatherings (\href{https://www.law.cornell.edu/supremecourt/text/24-1046}{Cornell LII}). Second, Haw. Rev. Stat. \S{}134-9.5(a)---dubbed the ``vampire rule'' by critics because folklore holds that a vampire cannot cross a threshold without an invitation---prohibited licensed carriers from bringing any firearm (even unloaded or inoperable) onto ``private property of another'' without ``express authorization'' from the owner, lessee, operator, manager, or authorized agent (\href{https://www.law.cornell.edu/supremecourt/text/24-1046}{Cornell LII}; \href{https://reason.com/volokh/2026/10/28/how-wolford-v-lopez-refines-second-amendment-doctrine/}{Volokh Conspiracy}). Under \S{}134-9.5(b), express authorization required either ``clear and conspicuous signage'' welcoming firearms or ``[u]nambiguous written or verbal authorization'' (\href{https://www.law.cornell.edu/supremecourt/text/24-1046}{Cornell LII}).

Act 52 flipped the longstanding Anglo-American common-law default rule, under which opening private property to the general public (such as gas stations, grocery stores, restaurants, and retail shops) creates an implied license to enter---including while lawfully carrying a firearm---unless the owner revokes consent (\href{https://www.supremecourt.gov/opinions/25pdf/24-1046_nmio.pdf}{Slip Opinion}). Four other states singled out in \textit{Bruen}---California, Maryland, New Jersey, and New York---adopted similar opt-in default rules (\href{https://www.scotusblog.com/2026/06/supreme-court-strikes-hawaii-gun-restriction/}{SCOTUSblog}). On June 23, 2023, three Maui County residents holding concealed-carry permits---Jason Wolford, Alison Wolford, and Atom Kasprzycki---joined by the Hawaii Firearms Coalition, sued Hawaii Attorney General Anne E. Lopez in federal court, challenging Act 52's sensitive-place bans and its private-property default rule under the Second and Fourteenth Amendments (\href{https://www.law.cornell.edu/supct/cert/24-1046}{Cornell LII Bulletin}; \href{https://www.everytown.org/wolford-v-lopez-second-amendment-scotus-oral-argument/}{Everytown}).

\section{Procedural History}
In August 2023, U.S. District Judge Leslie E. Kobayashi of the U.S. District Court for the District of Hawaii granted in part and denied in part plaintiffs' motion for a temporary restraining order and preliminary injunction (\textit{Wolford v. Lopez}, 686 F. Supp. 3d 1034 (D. Haw. 2023)) (\href{https://harvardlawreview.org/print/vol-140/wolford-v-lopez/}{Harvard Law Review}). Judge Kobayashi declined to enjoin \S{}134-9.5 as applied to private property closed to the public, but enjoined enforcement of the express-consent default rule as applied to private property open to the public, finding that Hawaii's historical analogues did not justify a blanket opt-in rule for public-facing businesses (\href{https://www.law.cornell.edu/supremecourt/text/24-1046}{Cornell LII}; \href{https://harvardlawreview.org/print/vol-140/wolford-v-lopez/}{Harvard Law Review}).

On September 6, 2024, a three-judge panel of the U.S. Court of Appeals for the Ninth Circuit, consolidating Hawaii's appeal with challenges to California's post-\textit{Bruen} statute, reversed the preliminary injunction as to Hawaii's \S{}134-9.5 while striking down California's stricter signage-only rule (\textit{Wolford v. Lopez}, 116 F.4th 959 (9th Cir. 2024)) (\href{https://cdn.ca9.uscourts.gov/datastore/opinions/2024/09/06/23-16164.pdf}{Ninth Circuit Opinion}; \href{https://www.oyez.org/cases/2025/24-1046}{Oyez}). The panel reasoned that colonial and Reconstruction-era statutes established a national tradition permitting states to arrange default consent rules for armed entry onto private property, so long as owners could consent verbally, in writing, or by signage (\href{https://harvardlawreview.org/print/vol-140/wolford-v-lopez/}{Harvard Law Review}). On January 15, 2025, the full Ninth Circuit denied rehearing en banc over dissents by Judge VanDyke (joined by five judges) and Judge Collins (joined by Judge Bress) (125 F.4th 1230 (9th Cir. 2025)) (\href{https://cdn.ca9.uscourts.gov/datastore/opinions/2025/01/15/23-16164.pdf}{En Banc Denial}; \href{https://www.law.cornell.edu/supremecourt/text/24-1046}{Cornell LII}).

Petitioners sought certiorari on April 1, 2025, and the Supreme Court granted review on October 3, 2025 (606 U.S. 1066), limited to whether Hawaii may prohibit licensed concealed-carry permit holders from carrying handguns on private property open to the public without the owner's express permission (\href{https://supreme.justia.com/cases/federal/us/609/24-1046/}{Justia}). The Court heard oral argument on January 20, 2026---with the United States arguing as amicus curiae in support of petitioners---and decided the case on June 25, 2026, reversing and remanding (609 U.S. 185, 146 S. Ct. 2032 (2026)) (\href{https://supreme.justia.com/cases/federal/us/609/24-1046/}{Justia}). On August 10, 2026, on remand, the Ninth Circuit affirmed the district court's preliminary injunction against \S{}134-9.5 (\href{https://cdn.ca9.uscourts.gov/datastore/opinions/2026/08/10/23-16164.pdf}{Ninth Circuit Remand Order}).

\section{Judicial Votes}
\textbf{Vote:} 6--3, reversing and remanding the Ninth Circuit (\href{https://www.law.cornell.edu/supremecourt/text/24-1046}{Cornell LII}).
\begin{itemize}
\item \textbf{Majority (6):} Justice Alito delivered the opinion of the Court, joined in full by Chief Justice Roberts and Justices Thomas, Gorsuch, Kavanaugh, and Barrett.
\item \textbf{Concurrence:} Justice Barrett filed a concurring opinion, joined by Justices Thomas and Gorsuch as to Part II--B (addressing 19th-century Black Codes).
\item \textbf{Dissents (3):} Justice Kagan filed a one-paragraph dissenting opinion. Justice Jackson filed a principal dissenting opinion, joined by Justice Sotomayor.
\end{itemize}

\section{Holding}
\textbf{Judgment:} Reversed and remanded. The Ninth Circuit's reversal of the preliminary injunction against Haw. Rev. Stat. \S{}134-9.5 was overturned (\href{https://www.law.cornell.edu/supremecourt/text/24-1046}{Cornell LII}).

\textbf{Rule Announced:} Under the Second and Fourteenth Amendments, a state may not prohibit licensed concealed-carry permit holders from carrying handguns onto private property held open to the public unless the property owner provides express, affirmative authorization. Such a state-imposed opt-in default rule falls within the plain text of the Second Amendment at \textit{Bruen} step one and lacks a ``relevantly similar'' historical analogue in the Nation's tradition of firearm regulation at \textit{Bruen} step two (\href{https://www.supremecourt.gov/opinions/25pdf/24-1046_nmio.pdf}{Slip Opinion}; \href{https://supreme.justia.com/cases/federal/us/609/24-1046/}{Justia}).

\textbf{Holding vs. Dicta:}
\begin{itemize}
\item \textit{Binding Holding:} Haw. Rev. Stat. \S{}134-9.5 violates the Second and Fourteenth Amendments as applied to private property open to the public; historical limits derived from property-law traditions belong at \textit{Bruen} step two rather than step one; local customs (such as Hawaii's ``spirit of Aloha'') cannot dilute the uniform national scope of the Second Amendment; and colonial anti-poaching statutes, an 1893 Oregon trespass law, and an 1865 Louisiana Black Code are not relevantly similar analogues supporting a broad opt-in carry default for commercial establishments (\href{https://www.law.cornell.edu/supremecourt/text/24-1046}{Cornell LII}).
\item \textit{Dicta and Open Questions:} The Court did not disturb a private property owner's independent right to exclude firearms from their premises (e.g., by posting ``No Guns'' signs or revoking consent), nor did it rule on private property \textit{closed} to the public (such as private residences) or on Hawaii's separate sensitive-place bans under \S{}134-9.1 (\href{https://www.law.cornell.edu/supremecourt/text/24-1046}{Cornell LII}; \href{https://www.scotusblog.com/2026/06/supreme-court-strikes-hawaii-gun-restriction/}{SCOTUSblog}). The Court also expressly left unresolved whether 1791 or 1868 is the controlling timeframe for Second Amendment incorporation analysis and declined to lay down a comprehensive standard governing when racially discriminatory historical statutes may be considered at step two (\href{https://www.law.cornell.edu/supremecourt/text/24-1046}{Cornell LII}; \href{https://harvardlawreview.org/print/vol-140/wolford-v-lopez/}{Harvard Law Review}).
\end{itemize}

\section{Analysis of Opinions}

\subsection{Majority Opinion (Justice Alito)}
Justice Alito anchored the Court's analysis in \textit{Heller}, \textit{McDonald}, \textit{Bruen}, and \textit{United States v. Rahimi}, 602 U.S. 680 (2024) (\href{https://www.law.cornell.edu/supremecourt/text/24-1046}{Cornell LII}). At \textit{Bruen} step one, Alito held that \S{}134-9.5 easily falls within the Second Amendment's plain text: petitioners are members of ``the people'' seeking to ``bear'' ``Arms'' (handguns) for self-defense. He rejected Hawaii's argument that historical property-law limits exclude armed entry onto private property at step one, insisting that historical regulations belong at step two. Turning to practical burdens, Alito illustrated the statute's impact through a hypothetical licensed carrier modeled on Jaime Caetano (\textit{Caetano v. Massachusetts}, 577 U.S. 411 (2016)), showing that running ordinary daily errands---stopping at a gas station, lunch restaurant, drug store, dry cleaner, and grocery store---could make a law-abiding permit holder ``a criminal at least six times over'' because indifferent or image-conscious proprietors will rarely post ``Guns Welcome'' signs (\href{https://www.law.cornell.edu/supremecourt/text/24-1046}{Cornell LII}).

At \textit{Bruen} step two, Alito dismantled Hawaii's three historical defenses. First, invoking \textit{McDonald}, he rejected Hawaii's reliance on its unique pre-statehood weapon restrictions and ``the spirit of Aloha'' (\textit{State v. Wilson}, 154 Haw. 8 (2024)), holding that fundamental Bill of Rights guarantees have a uniform national meaning across all 50 states (\href{https://www.law.cornell.edu/supremecourt/text/24-1046}{Cornell LII}). Second, Alito distinguished Hawaii's colonial and founding-era statutes---1721 Pennsylvania, 1722 New Jersey, 1728 Maryland, 1763 New York, and 1771 New Jersey---as anti-poaching laws targeting unauthorized hunting on rural or enclosed plantations where game could be found, designed to prevent game theft, stray gunfire, and disturbance of livestock, rather than peaceful concealed carry in retail businesses (\href{https://www.supremecourt.gov/opinions/25pdf/24-1046_nmio.pdf}{Slip Opinion}). Third, he dismissed an 1893 Oregon enclosed-land statute as an ambiguous late-19th-century outlier, and condemned Hawaii's reliance on an 1865 Louisiana statute as a ``tainted artifact'' from the post-Civil War Black Codes enacted to disarm freed Black citizens---a statute neither widespread nor accepted and antithetical to the Fourteenth Amendment's original understanding (\href{https://www.law.cornell.edu/supremecourt/text/24-1046}{Cornell LII}).

\subsection{Concurring Opinion (Justice Barrett)}
Justice Barrett joined the Court's opinion in full and wrote separately---joined by Justices Thomas and Gorsuch as to Part II--B---to sharpen the step-one and step-two methodology (\href{https://www.law.cornell.edu/supremecourt/text/24-1046}{Cornell LII}). In Part I, Barrett responded directly to Justice Jackson's property-law framing: while private owners may freely exclude firearms, the Second Amendment binds \textit{state action}, and a state property law targeting arms-bearing conduct triggers constitutional scrutiny just as a state default rule requiring express permission to wear a hijab into public-facing businesses would trigger First Amendment scrutiny. In Part II--A, Barrett agreed with Justice Jackson that defining the historical principle strictly as ``poaching'' would demand an improper ``historical twin,'' instead extracting a broader principle from the 18th-century statutes: states may adopt focused default rules in specific locations prone to particular abuses of firearms. Hawaii's rule still failed because it targeted no localized abuse, instead resting on generalized public disapproval of guns across all commercial property. In Part II--B (joined by Thomas and Gorsuch), Barrett documented the history of Louisiana's 1865 law and parallel 1866 Texas and Florida Black Codes, concluding that even setting aside whether morally odious laws are admissible evidence, their ``why''---subordinating freedmen and forcing them into plantation labor---was not relevantly similar to Hawaii's justification (\href{https://www.law.cornell.edu/supremecourt/text/24-1046}{Cornell LII}; \href{https://harvardlawreview.org/print/vol-140/wolford-v-lopez/}{Harvard Law Review}).

\subsection{Dissenting Opinion (Justice Kagan)}
In a crisp one-paragraph dissent, Justice Kagan argued that Hawaii's statute should be upheld solely at \textit{Bruen} step two based on the colonial and founding-era laws analyzed in Part III of Justice Jackson's dissent (\href{https://www.law.cornell.edu/supremecourt/text/24-1046}{Cornell LII}). In Kagan's view, the ``how'' was identical---setting a no-carry default that any private landowner remains free to reverse---and the ``why'' was sufficiently analogous under \textit{Rahimi}'s principles-based test because both the founding-era laws and Act 52 address the ``abuses, damages and inconveniences'' associated with carrying firearms on another person's property. Because that sufficed to uphold the law, Kagan expressly declined to address \textit{Bruen} step one or the evidentiary status of Louisiana's 1865 Black Code (\href{https://www.law.cornell.edu/supremecourt/text/24-1046}{Cornell LII}).

\subsection{Dissenting Opinion (Justice Jackson)}
Justice Jackson, joined by Justice Sotomayor, framed the dispute as ``about property rights, not gun rights'' (\href{https://www.law.cornell.edu/supremecourt/text/24-1046}{Cornell LII}). At \textit{Bruen} step one, Jackson argued that because the Second Amendment codified a pre-existing right subject to the fundamental property-law right to exclude, and because all parties conceded there is no constitutional right to carry a firearm onto private property without the owner's consent, state rules specifying the \textit{form} of that consent (express versus implied) do not implicate the Second Amendment at all. She accused the majority of ``mov[ing] the goal posts'' by stripping history out of step one and turning it into a ``one-way ratchet'' that renders every gun regulation presumptively unconstitutional (\href{https://www.law.cornell.edu/supremecourt/text/24-1046}{Cornell LII}).

At step two, Jackson contended that colonial and early state laws---especially New Jersey's 1771 statute covering taxed lands (which historians noted included taverns and shops) and statutes covering ``premises''---demonstrated an established tradition of opt-in defaults protecting property owners' right to exclude and preventing harms from unauthorized armed entry. Demanding a closer match to poaching, she argued, violated \textit{Rahimi}'s warning against trapping the law ``in amber.'' Finally, Jackson sharply criticized the majority's categorical exclusion of Louisiana's 1865 Black Code: while condemning the Black Codes as ``ugly,'' ``racist,'' and ``deplorable,'' she argued that if the Court insists on tethering Second Amendment doctrine to history, it must distinguish Fourteenth Amendment equal-protection violations from Second Amendment scope and confront America's entire historical record ``warts and all,'' rather than selectively excising uncomfortable analogues to invalidate modern gun laws (\href{https://www.law.cornell.edu/supremecourt/text/24-1046}{Cornell LII}; \href{https://www.scotusblog.com/2026/06/supreme-court-strikes-hawaii-gun-restriction/}{SCOTUSblog}).

\section{Examples: Future Applications}
\begin{enumerate}
\item \textbf{Same-Side (Unconstitutional):} \textit{California's Signage-Only Opt-In Default Rule.} A state statute bars licensed concealed-carry permit holders from entering any private commercial establishment open to the public unless the business posts a state-prescribed, standardized ``Firearms Permitted'' placard at every entrance. \textit{Why:} This law is even more restrictive than Hawaii's \S{}134-9.5 because it eliminates verbal and individualized written consent, directly violating \textit{Wolford}'s holding that states may not flip the common-law implied-license default for public-facing private property (\href{https://www.law.cornell.edu/supremecourt/text/24-1046}{Cornell LII}; \href{https://www.scotusblog.com/2026/06/supreme-court-strikes-hawaii-gun-restriction/}{SCOTUSblog}).

\item \textbf{Same-Side (Unconstitutional):} \textit{Municipal Opt-In Rule for Privately Owned Parking Lots.} A city ordinance makes it a misdemeanor for licensed gun owners to bring a firearm---even concealed on their person or locked in a vehicle---onto any privately owned retail parking lot open to customers unless the lot owner files an annual affirmative opt-in registration with the city clerk. \textit{Why:} Like Hawaii's law, this flips the common-law implied license on commercial property open to the public and functionally traps permit holders in daily errands without any relevantly similar historical analogue (\href{https://www.supremecourt.gov/opinions/25pdf/24-1046_nmio.pdf}{Slip Opinion}).

\item \textbf{Opposite-Side (Constitutional):} \textit{Private Business ``No Firearms'' Signage Enforced via Criminal Trespass.} A statewide grocery chain posts ``No Firearms Allowed'' signs at its store entrances, and the state enforces its general criminal trespass statute against a licensed permit holder who knowingly ignores the sign and refuses to leave. \textit{Why:} Both the majority and Justice Barrett's concurrence reaffirmed that the Second Amendment restricts only state action and leaves private property owners completely free to revoke implied consent and exclude armed visitors using state trespass remedies (\href{https://www.law.cornell.edu/supremecourt/text/24-1046}{Cornell LII}).

\item \textbf{Opposite-Side (Constitutional):} \textit{Express-Consent Default for Private Residences and Non-Public Farms.} A state statute prohibits carrying firearms onto private residences, fenced residential yards, and non-public agricultural or timber lands without the owner's express verbal or written permission. \textit{Why:} Such property is \textit{closed} to the general public, so there is no common-law implied license to enter, and the law closely parallels the colonial and founding-era enclosed-land and plantation trespass statutes discussed in \textit{Wolford} (\href{https://www.law.cornell.edu/supremecourt/text/24-1046}{Cornell LII}).

\item \textbf{Fence-Sitter (Close Call):} \textit{Targeted Opt-In Default for High-Risk Commercial Venues.} Following a wave of late-night shootings at 24-hour nightclubs and cannabis dispensaries, a state enacts an opt-in default requiring express owner consent to carry firearms into those two narrow categories of private businesses (which are open to the public, though not covered by the state's flat sensitive-place bans). \textit{Why:} Under Justice Alito's opinion, the law still flips the common-law default on public-facing private property and lacks an exact colonial retail analogue; yet under Justice Barrett's concurrence (Part II--A) and \textit{Rahimi}, the state could argue it is a ``focused regulation'' targeting specific places prone to documented, localized abuses of firearms rather than generalized disapproval of public carry (\href{https://www.law.cornell.edu/supremecourt/text/24-1046}{Cornell LII}).
\end{enumerate}

\section{Critique}
\textbf{Originalist and Textualist Perspectives:} Supporters of the decision---including scholars such as Joseph G.S. Greenlee and George A. Mocsary---praised \textit{Wolford} for blocking post-\textit{Bruen} ``evasion'' regimes that would convert an enumerated constitutional right into a government-conditioned permission slip across everyday commercial life (\href{https://firearmsresearchcenter.org/working_papers/vampire-rules-monarchical-legacies-and-black-codes-wolford-v-lopez-and-the-limits-of-post-bruen-evasion/}{Firearms Research Center}; \href{https://reason.com/volokh/2026/10/28/how-wolford-v-lopez-refines-second-amendment-doctrine/}{Volokh Conspiracy}). From this vantage point, distinguishing private owner opt-out authority from a state-mandated opt-in baseline preserves both property rights and public carry. Yet even sympathetic originalists noted methodological tensions: Justice Alito's narrow framing of the 18th-century statutes as ``anti-poaching'' laws risked reverting toward the ``historical twin'' rigidity that \textit{Rahimi} disavowed---a tension Justice Barrett explicitly sought to moderate in Part II--A of her concurrence (\href{https://www.law.cornell.edu/supremecourt/text/24-1046}{Cornell LII}).

\textbf{Progressive, Property-Law, and Doctrinal Perspectives:} Critics, including the \textit{Harvard Law Review} and property-law scholar Maureen E. Brady, argued that \textit{Wolford} exposes deep indeterminacy in \textit{Bruen}'s history-and-tradition framework (\href{https://harvardlawreview.org/print/vol-140/wolford-v-lopez/}{Harvard Law Review}; \href{https://review.law.stanford.edu/online/property-v-guns-the-level-of-generality-problem-in-wolford/}{Stanford Law Review Online}). First, Brady highlighted the ``level-of-generality problem'': courts can manipulate outcomes by characterizing colonial statutes narrowly as ``anti-poaching'' rules or broadly as ``armed-trespass default'' rules. Second, the \textit{Harvard Law Review} noted that Alito's exclusion of Louisiana's 1865 Black Code creates an unguided ``\textit{Bruen} step 1.5,'' forcing lower courts to make ad hoc moral judgments about which racially or religiously discriminatory historical gun laws are too ``tainted'' to count as usable history (\href{https://harvardlawreview.org/print/vol-140/wolford-v-lopez/}{Harvard Law Review}). Gun-safety advocates, such as Everytown, further argued that the decision overrides state federalism and burdens small businesses that do not want firearms on their premises but fear backlash from posting ``No Guns'' signs (\href{https://www.everytown.org/wolford-v-lopez-second-amendment-scotus-decision-explained/}{Everytown}).

\textbf{Underlying Competing Values:} \textit{Wolford} reflects an enduring clash among three constitutional values: (1) ensuring that an individual's enumerated right to armed self-defense remains practically exercisable in ordinary daily life; (2) preserving state sovereignty to define background property-law default rules and respond to local public-safety norms; and (3) maintaining methodological integrity in historical jurisprudence when the Nation's past includes racially discriminatory regulations that courts are understandably reluctant to validate.

\section{Key Quotations}

\begin{quote}
``This regime hobbles what the Second Amendment protects: the right of Americans to carry arms for self-defense as they go about their daily lives. We hold that the law is unconstitutional.''
\end{quote}
\noindent --- Justice Alito, Opinion of the Court (\href{https://www.law.cornell.edu/supremecourt/text/24-1046}{Cornell LII}; \href{https://www.courtlistener.com/opinion/10880241/wolford-v-lopez/}{CourtListener}).

\begin{quote}
``It cannot give way to `the spirit of Aloha' in Hawaii, contra, \textit{State v. Wilson}, 154 Haw. 8, 27, 543 P. 3d 440, 459 (2024), any more than it can yield to the spirit of the Big Apple (\textit{Bruen}) or the Windy City (\textit{McDonald}).''
\end{quote}
\noindent --- Justice Alito, Opinion of the Court (\href{https://www.law.cornell.edu/supremecourt/text/24-1046}{Cornell LII}).

\begin{quote}
``But the Second Amendment does not apply to private parties. It does apply to the States. See U. S. Const., Amdt. 14, \S{}1. And when a State enacts a property law that regulates arms-bearing conduct, that law implicates the Second Amendment.''
\end{quote}
\noindent --- Justice Barrett, concurring (\href{https://www.law.cornell.edu/supremecourt/text/24-1046}{Cornell LII}).

\begin{quote}
``The `how' is identical: The new law, just like the old ones, sets a default rule against gun carry that a private landowner may reverse. The `why' is sufficiently close. Both sets of laws respond to the dangers and harms that someone with a gun can cause on another person's property.''
\end{quote}
\noindent --- Justice Kagan, dissenting (\href{https://www.law.cornell.edu/supremecourt/text/24-1046}{Cornell LII}).

\begin{quote}
``With this decision, the Court has now manipulated \textit{Bruen} into a free-for-all that lets the Judiciary thwart the will of legislatures by privileging access to firearms above all else. Today's decision makes one thing clear: The Court's objective is protecting guns, not consistently preserving any principle of law.''
\end{quote}
\noindent --- Justice Jackson, dissenting (\href{https://www.law.cornell.edu/supremecourt/text/24-1046}{Cornell LII}; \href{https://supreme.justia.com/cases/federal/us/609/24-1046/}{Justia}).

\end{document}
